\documentclass[11pt]{article}
\usepackage[margin=0.85in]{geometry}
\usepackage{lmodern}
\usepackage{microtype}
\usepackage{enumitem}
\usepackage{listings}
\usepackage[T1]{fontenc}

\lstset{
  basicstyle=\ttfamily\small,
  columns=fullflexible,
  keepspaces=true,
  showstringspaces=false,
  breaklines=true,
  frame=single,
  xleftmargin=0.4em,
  xrightmargin=0.4em,
  aboveskip=0.5em,
  belowskip=0.5em
}

\pagestyle{empty}
\setlength{\parskip}{0.32em}
\setlength{\parindent}{0pt}
\setlist[enumerate,1]{leftmargin=1.4em, itemsep=0.65em, topsep=0.25em}
\setlist[enumerate,2]{leftmargin=1.5em, itemsep=0.28em, topsep=0.2em, label=\alph*.}
\setlist[itemize]{leftmargin=1.3em, itemsep=0.12em, topsep=0.15em}

\begin{document}

\begin{center}
{\LARGE\bfseries Final exam --- Sample 2}\\[0.35em]
DATA 412/612
\end{center}

\vspace{0.2em}

Practice exam, questions only. Cumulative through December 9
(weeks 1--14). The newest material is lists, loops, and
\texttt{purrr} (\texttt{pmap}, \texttt{keep}, \texttt{discard}).
Earlier weeks stay in scope, including a report chunk.

Put your answers in an R Markdown document and keep the code visible
except where a chunk option below says to hide it.
You may use your notes, the R help files, and the internet.
Do not discuss the exam with anyone else.
Load the tidyverse.

\begin{enumerate}
\item \textbf{A loop and a map} (8 pts).
\texttt{bases} is \texttt{c(2, 3, 4)}.
With a pre-allocated \texttt{for} loop, fill a numeric vector with
the square of each entry.
Then do the same job with \texttt{map\_dbl()}.
The two results should match.

\item \textbf{\texttt{pmap}} (6 pts).
Using the same means and standard deviations as
\texttt{mu = c(0, 10)} and \texttt{sigma = c(1, 2)}, draw 4 normal
values from each pair with \texttt{pmap()} and a list of arguments.
Set the seed to \texttt{612} first.

\item \textbf{\texttt{keep} and \texttt{discard}} (6 pts).
From \texttt{mpg}, \texttt{keep} the numeric columns and report how
many there are.
Then \texttt{discard} every column that is not a character column,
and print the names that remain.

\item \textbf{Tidy, then test} (12 pts).
Reshape \texttt{wide} to a \texttt{year} column and a \texttt{sales}
column. \texttt{year} is numeric. Drop the missing sales.
Turn \texttt{store} into a factor ordered by total sales, largest
first.
Run a two-sample \texttt{t.test()} of \texttt{sales} for store
\texttt{A} against store \texttt{B}. State the null hypothesis in
one sentence.

\begin{lstlisting}
wide <- tibble(
  store = c("A", "B", "A"),
  y2022 = c(5, 9, NA),
  y2023 = c(7, 4, 6)
)
\end{lstlisting}

\item \textbf{The report around that test} (8 pts).
Put the test in a chunk labeled \texttt{store-test} with
\texttt{echo = TRUE}.
Put the reshape in a setup-style chunk with \texttt{include = FALSE}.
Send the store totals to \texttt{knitr::kable()} with a caption.
In YAML, set a bibliography file \texttt{refs.bib}, and in the text
cite \texttt{wickham2017r4ds} in parentheses.
\end{enumerate}

\end{document}
